% Shared body for Paper B (agent-auth SoK / position, work-in-progress).
% Class-agnostic; the per-venue main.tex sets class/title/authors.
\providecommand{\gate}[1]{\textbf{[PENDING: #1]}}

\begin{abstract}
Autonomous agents increasingly act on a user's behalf carrying OAuth/MCP bearer
tokens, yet the authorization layer around them is new and largely unstudied. We
systematize the authorization threat model for agent and Model Context Protocol
(MCP) deployments, mapping each threat to its normative source (OAuth RFCs
8693/8707/9728, the MCP authorization specification, OWASP LLM \& Agentic Top~10s)
and to a concrete, checkable invariant. We contribute a cited taxonomy of 15
agent-domain threats, a gap analysis of what existing OAuth machinery does and does
not cover for autonomous callers, and an open-source static analyzer that flags
audience-binding, delegation, scope, and expiry violations on a single token.
\gate{a measurement of prevalence across public MCP deployments; this draft is
work-in-progress / position and the empirical section is not yet complete.}
\end{abstract}

\section{Introduction}
When an agent calls a tool or an MCP server on a user's behalf, a bearer token
travels with the request. The failure modes are OAuth failure modes with sharper
consequences because the caller is autonomous: a token \emph{not bound to the
resource} it is presented to can be replayed or passed through (a confused deputy);
an on-behalf-of token with no verifiable delegation chain makes ``who authorised
this?'' unanswerable; an omnibus scope (\texttt{admin:*}) hands an injected or
misbehaving agent the keys to everything; and a non-expiring credential never stops
being useful to an attacker. These map directly onto the MCP authorization
specification, OAuth token exchange, and the OWASP LLM/Agentic Top~10s---yet no
systematic, source-grounded treatment exists. Being early and rigorous here is the
contribution.

\section{Threat systematization}
We organize the agent-authorization threat space into categories---audience
binding, delegation, scope, expiry, agent identity, and agency---and, for each
threat, give its definition, normative source, the invariant that must hold, and
the failure it prevents. Representative entries with checkable invariants:
\begin{itemize}
  \item \textbf{Unbound audience (MCP-02).} A token carrying no audience is not
  bound to any resource and is replayable (RFC~8707/9728).
  \item \textbf{Audience mismatch (MCP-01).} A token whose audience omits the
  presented resource can be passed through to another server (RFC~8707).
  \item \textbf{Missing delegation (DEL-01).} An on-behalf-of token without a
  verifiable \texttt{act} chain has an unauditable delegation (RFC~8693).
  \item \textbf{Over-broad scope (SCOPE-01).} Omnibus grants violate least
  privilege (OWASP Agentic).
  \item \textbf{Non-expiring credential (SCOPE-02).} No \texttt{exp} / excessive
  lifetime (OAuth security BCP).
\end{itemize}
The full taxonomy (this work) spans 15 agent-domain threats, all cited.

\section{Gap analysis}
We contrast the classic OAuth/OIDC threat model with the autonomous-agent setting:
which threats existing machinery (resource indicators, token exchange, protected-
resource metadata) covers by default in common SDKs, and which are left to the
deployer and therefore likely missed (e.g., confused-deputy via token passthrough;
delegation chains that token exchange models but MCP deployments rarely verify).

\section{Instrument}
We provide an open-source static analyzer (\texttt{keyway agent inspect}) that
decodes a token and emits cited findings against a supplied policy: audience
binding, on-behalf-of delegation, scope minimisation, and expiry. It is stateless
and runs on a single token---in CI or client-side---so nothing leaves the caller.
Of the 15-threat agent taxonomy, the analyzer currently covers 6 (40\%); every gap
is named.

\section{Measurement (work in progress)}
\gate{measure the prevalence of the checkable weaknesses above across a corpus of
public MCP servers and reference agent-token deployments, with a hand-validated
sample. This is the empirical contribution and is not yet complete; the present
submission is positioned as SoK / position / work-in-progress accordingly.}

\section{Threats to validity and ethics}
The MCP specification is fast-moving; the systematization is date-stamped and
coverage is a snapshot. Public MCP servers skew to early adopters. We probe no
third-party live agents and use no real user tokens; identifiable high-severity
findings follow responsible disclosure.

\section{Conclusion}
Agent authorization is an under-studied, high-stakes surface. We contribute a
source-grounded threat systematization, a gap analysis, and an open-source
instrument, as the foundation for a forthcoming measurement of the space.
